\section{Introduction}
\label{sec:introduction}

% Fig. 1: GaitNet pipeline overview. Layout follows the FADA method
% overview (Xie et al.): training on the training controller on the left;
% data collection, IDM refit and deployment on the new controller on the
% right. The deployment panel shows GaitNet's architecture and the
% execution layer (IDM for where, timing for when) as separate blocks;
% domain bands below. Styles and colours: fig_style.tex. Architecture
% follows gaitnet-docs/methodology.md 5-6 (detailed version:
% fig_architecture.tex); step timing follows Sec. III-E.
\begin{figure*}[t]
\centering
\resizebox{\textwidth}{!}{%
\begin{tikzpicture}[gfig]
% ---------------------------------------------------------------- panels
\draw[region={rgS}{rgSd}] (0,0)        rectangle (6.6,-2.75);    % (a)
\draw[region={rgS}{rgSd}] (0,-2.95)    rectangle (6.6,-6.35);    % (b)
\draw[region={rgG}{rgGd}] (6.85,0)     rectangle (12.25,-2.75);  % (c)
\draw[region={rgG}{rgGd}] (12.5,0)     rectangle (19.6,-2.75);   % (d)
\draw[region={rgG}{rgGd}] (6.85,-2.95) rectangle (19.6,-6.35);   % (e)
\node[ptitle] at (0.2,-0.2)   {(a) IDM training};
\node[ptitle] at (0.2,-3.15)  {(b) GaitNet training};
\node[ptitle] at (7.05,-0.2)  {(c) Data collection on OCS2};
\node[ptitle] at (12.7,-0.2)  {(d) IDM refit for OCS2};
\node[ptitle] at (7.05,-3.15) {(e) Deployment};
% domain bands
\draw[band={rgS}{rgSd}] (0,-6.6)    rectangle (6.6,-7.15);
\draw[band={rgG}{rgGd}] (6.85,-6.6) rectangle (19.6,-7.15);
\node[font=\sffamily\footnotesize\bfseries, text=ink] at (3.3,-6.875)
  {Isaac Lab, training controller};
\node[font=\sffamily\footnotesize\bfseries, text=ink] at (13.225,-6.875)
  {OCS2, new controller};

% ---------------------------------------------------------------- (a)
% IDM fitted by hindsight relabelling: input the landing o, target the
% command u that produced it.
\def\ya{-1.6}
\node[mod=acC, minimum width=16mm, minimum height=14mm] (ia) at (3.0,\ya)
  {IDM\\[1pt]$I^{\text{train}}_\theta$ \hot};
\node[font=\sffamily\footnotesize, anchor=east] (ao) at (1.55,{\ya+0.35}) {landing $\mathbf o$};
\node[font=\sffamily\footnotesize, anchor=east] (as) at (1.55,{\ya-0.35}) {$\mathbf s,\ell,T$};
\draw[arr] (ao.east) -- (ao.east -| ia.west);
\draw[arr] (as.east) -- (as.east -| ia.west);
\node (ua) at (4.45,\ya) {$\hat{\mathbf u}$};
\draw[arr] (ia) -- (ua);
\node[mod=acT, minimum width=10mm, text=acT] (ma) at (5.55,\ya) {MSE};
\draw[arr] (ua) -- (ma);
\node[font=\sffamily\footnotesize] (ta) at (5.55,-0.55) {command $\mathbf u$};
\draw[arr] (ta) -- (ma);
\node[sub] at (5.2,-2.3) {relabelled in hindsight:\\input $\mathbf o$, target $\mathbf u$};

% ---------------------------------------------------------------- (b)
% PPO with the training controller's IDM fixed between GaitNet and the
% (randomized) training controller.
\def\yb{-4.55}
\node[mod=acB, minimum width=15mm, minimum height=10.5mm] (gb) at (1.0,\yb) {GaitNet\\[1pt]\hot};
\node[mod=acC, minimum width=15mm, minimum height=10.5mm] (ib) at (3.0,\yb) {IDM\\[1pt]$I^{\text{train}}_\theta$ \frz};
\node[mod=muted, minimum width=19mm, minimum height=10.5mm] (cb) at (5.2,\yb) {training\\controller};
\draw[arr] (gb) -- node[above, alab] {$\mathbf u^\star$} (ib);
\draw[arr] (ib) -- node[above, alab] {$\mathbf u$} (cb);
\draw[tarr] (cb.south) -- (5.2,-5.75) -- node[above, tlab, text=acT]
  {reward, PPO update} (1.0,-5.75) -- (gb.south);
\draw[->, line width=0.7pt, draw=acC, dash pattern=on 2pt off 1.5pt]
  (ia.south) -- node[right, tlab, text=muted, pos=0.82] {fixed} (ib.north);

% ---------------------------------------------------------------- (c)
% A policy walks on OCS2 and every executed step is logged as (o, u).
\def\yc{-1.3}
\node[mod=ink, minimum width=16mm, minimum height=10mm, rounded corners=8pt] (pc) at (8.1,\yc) {};
\node[font=\sffamily\footnotesize] at (8.1,{\yc+0.27}) {policy};
\pic[scale=0.7] at (8.1,{\yc-0.13}) {mlp=acB};
\node[mod=ink, minimum width=17mm, minimum height=10mm, rounded corners=8pt] (oc) at (10.95,\yc)
  {OCS2\\[-1pt]{\scriptsize NMPC + WBC}};
\draw[arr] ([yshift=1.6mm]pc.east) -- node[above, alab] {$\mathbf u$} ([yshift=1.6mm]oc.west);
\draw[arr] ([yshift=-1.6mm]oc.west) -- node[below, alab] {$\mathbf o$} ([yshift=-1.6mm]pc.east);
\node[font=\sffamily\scriptsize, text=ink, anchor=west, inner sep=0pt] at (pc.west |- 0,-2.3) {logged steps};
\node[tok={tkS}{acB!45}, minimum width=6.5mm] at (9.2,-2.3)  {$\mathbf o_1$};
\node[tok={white}{acC},  minimum width=6.5mm] at (9.85,-2.3) {$\mathbf u_1$};
\node[tok={tkS}{acB!45}, minimum width=6.5mm] at (10.6,-2.3) {$\mathbf o_2$};
\node[tok={white}{acC},  minimum width=6.5mm] at (11.25,-2.3) {$\mathbf u_2$};
\node[inner sep=1pt] (dots) at (11.85,-2.3) {$\cdots$};

% ---------------------------------------------------------------- (d)
% The logged landing o is the IDM's input and the logged command u its
% target (hindsight relabelling, as in (a)); GaitNet is not retrained.
\def\yd{-1.3}
\node[mod=acC, minimum width=17mm, minimum height=10mm] (id) at (14.4,\yd)
  {IDM \hot\\[-1pt]{\scriptsize $\mathbf o + \Delta^{\text{OCS2}}_\theta$}};
\node[mod=acT, minimum width=10mm, text=acT] (md) at (16.55,\yd) {MSE};
\draw[arr] (id) -- node[above, alab] {$\hat{\mathbf u}$} (md);
\coordinate (jd) at (12.95,-2.3);
\draw[line width=0.7pt, draw=ink] (dots.east) -- (jd);
\fill[ink] (jd) circle (0.05);
\draw[arr] (jd) |- node[pos=0.75, above, alab] {$\mathbf o$} (id.west);
\draw[tarr] (jd) -| node[pos=0.25, above, tlab, text=acT] {target $\mathbf u$} (md.south);
\draw[muted, line width=0.6pt, dash pattern=on 2pt off 2pt] (17.6,-0.65) -- (17.6,-2.5);
\node[mod=acB, minimum width=14mm, minimum height=10mm] (gd) at (18.6,\yd) {GaitNet\\[1pt]\frz};
\node[sub, anchor=north] at (18.6,-1.88) {unchanged};

% ---------------------------------------------------------------- (e)
% GaitNet (frozen) as in Sec. III-B: context z = f_s(x_g) + W_h f_h(h);
% tokens W_t[z; e_i] with e_i = f_c(q_i) and a learned e_0 for no-step;
% candidate mixer; heads (alpha_i, T_i); masked argmax. The execution
% layer (Sec. III-C, III-E) sets where (IDM) and when (timing) the step
% is executed on OCS2.
\def\yA{-4.25} \def\yB{-4.78} \def\yC{-5.31} \def\yD{-5.84}
\def\yTop{-4.03} \def\yBot{-6.06}   % top of row A and bottom of row D tokens
\draw[fill=rgB, draw=acB!70, line width=0.6pt, rounded corners=5pt] (7.05,-3.5) rectangle (15.45,-6.2);
\node[font=\sffamily\footnotesize\bfseries, text=acB, anchor=north west, inner sep=0pt] at (7.275,-3.6) {GaitNet \frz};
% inputs
\node[tok={tkS}{acB!45}, minimum height=4.4mm] (xe) at (7.6,\yA) {$\mathbf x_g$};
\node[tok={tkS}{acB!45}, minimum height=4.4mm] (he) at (7.6,\yB) {$\mathbf h$};
\node[tok={tkC}{tkCd},   minimum height=4.4mm] (qe) at (7.6,\yC) {$\mathbf q_i$};
\node[tok={tkN}{rgAd},   minimum height=4.4mm] (ne) at (7.6,\yD) {$\varnothing$};
% context z = f_s(x_g) + W_h f_h(h)
\node[mod=acB, minimum width=7mm, minimum height=4.4mm, inner sep=1pt] (fse) at (8.6,\yA) {$f_s$};
\node[mod=acB, minimum width=7mm, minimum height=4.4mm, inner sep=1pt] (fhe) at (8.6,\yB) {$f_h$};
\node[mod=acB, minimum width=7mm, minimum height=4.4mm, inner sep=1pt] (whe) at (9.6,\yB) {$W_h$};
\node[draw=acB, circle, line width=0.7pt, fill=white, inner sep=0.2pt] (pe) at (10.3,\yA) {$+$};
\node[tok={tkS}{acB!45}, minimum height=4.4mm] (ze) at (11.1,\yA) {$\mathbf z$};
\draw[arr] (xe) -- (fse);  \draw[arr] (he) -- (fhe);
\draw[arr] (fse) -- (pe);  \draw[arr] (fhe) -- (whe);
\draw[arr] (whe.east) -| (pe.south);
\draw[arr] (pe) -- (ze);
% candidate embeddings and tokens [z; e_i]; W_t spans rows C-D
\node[mod=acB, minimum width=7mm, minimum height=4.4mm, inner sep=1pt] (fce) at (8.6,\yC) {$f_c$};
\node[tok={tkN}{acB!45}, minimum width=7mm, minimum height=4.4mm] (e0e) at (8.6,\yD) {$\mathbf e_\varnothing$};
\draw[arr] (qe) -- (fce);  \draw[arr] (ne) -- (e0e);
\node[mod=acB, minimum width=9mm, minimum height=9.7mm, inner sep=1pt] (wte) at (11.1,{(\yC+\yD)/2})
  {$W_t$\\[-2pt]{\scriptsize $[\mathbf z;\mathbf e_i]$}};
\draw[arr] (fce) -- node[above, alab, font=\scriptsize, pos=0.7] {$\mathbf e_i$} (fce -| wte.west);
\draw[arr] (e0e) -- (e0e -| wte.west);
\draw[arr] (ze) -- (wte);
% candidate mixer spans rows B-D
\node[mod=acB, minimum width=20mm, minimum height=15mm] (mxe) at (12.85,{(\yB+\yD)/2}) {};
\node[font=\sffamily\scriptsize\bfseries, text=acB] at (12.85,{\yB-0.05}) {candidate mixer};
\pic[scale=0.8] at (12.85,\yC) {attn=acB};
\node[sub] at (12.85,{\yD+0.05}) {self-attn $\times$2};
\draw[arr] (wte) -- (wte -| mxe.west);
% heads (alpha_i, T_i) on the row-D baseline, masked argmax up to row A
\begin{scope}[shift={(14.2,\yBot)}]
  \foreach \x/\hgt/\c in {0/0.35/acB!35, 0.15/0.6/acB!35, 0.3/0.85/tkCd, 0.45/0.45/acB!35, 0.6/0.25/acB!35}
    \fill[\c] (\x,0) rectangle ++(0.11,\hgt);
  \draw[muted, line width=0.4pt] (-0.05,0) -- (0.76,0);
\end{scope}
\draw[arr] (mxe.east |- 0,\yC) -- (14.12,\yC);
\node[tok={tkC}{tkCd}, minimum width=14mm, minimum height=4.4mm] (se) at (14.555,\yA) {$(\ell,\mathbf u^\star\!,T)$};
\node[sub, anchor=south] at (14.555,{\yTop+0.02}) {masked argmax};
\draw[arr] (14.555,{\yBot+0.9}) -- (se);
% execution layer: where (IDM, row A) and when (timing, row C)
\draw[fill=white, draw=acC!70, line width=0.6pt, rounded corners=5pt] (15.7,-3.5) rectangle (18.25,-6.2);
\node[font=\sffamily\footnotesize\bfseries, text=acC, anchor=north west, inner sep=0pt] at (15.95,-3.6) {Execution layer};
\coordinate (je) at (15.95,\yA);
\draw[line width=0.7pt, draw=ink] (se.east) -- (je);
\fill[ink] (je) circle (0.05);
\node[mod=acC, minimum width=14mm, minimum height=6mm] (ie) at (17.15,\yA) {IDM \frz};
\node[sub, anchor=north] at (17.15,{\yA-0.33}) {where the foot lands};
\node[mod=acC, minimum width=14mm, minimum height=6mm] (tie) at (17.15,\yC) {timing};
\node[sub, anchor=north] at (17.15,{\yC-0.33}) {when the step starts};
\draw[arr] (je) -- (ie);
\draw[arr] (je) |- (tie.west);
% OCS2 receives the command and the start time; centred on row B
\node[mod=ink, minimum width=9.5mm, minimum height=15mm] (oe) at (18.95,\yB) {OCS2};
\draw[arr] (ie.east) -- node[pos=0.3, above, alab] {$\mathbf u$} (ie.east -| oe.west);
\draw[arr] (tie.east) -- (tie.east -| oe.west);
\end{tikzpicture}%
}
\caption{GaitNet pipeline. (a)~An inverse footstep model (IDM) for the
  training controller is fitted on its logged steps: given where a foot
  landed, it predicts the command that put it there. (b)~GaitNet is
  trained with PPO, with this IDM fixed between GaitNet and the
  controller. (c)~On a new controller, OCS2, steps are logged
  again\pending{ (policy to be added)}, and (d)~a new IDM is fitted on
  them while GaitNet stays frozen. (e)~At deployment, GaitNet encodes the
  state $\mathbf x_g$ and its recent motion $\mathbf h$ into a context
  $\mathbf z$, scores each safe candidate $\mathbf q_i$ and the no-step
  option $\varnothing$ in a candidate mixer, and picks one step
  $(\ell,\mathbf u^\star\!,T)$ (Fig.~\ref{fig:architecture}). The
  execution layer then sets where and when OCS2 executes it: the OCS2 IDM
  turns $\mathbf u^\star$ into a command $\mathbf u$, and the step is
  written into OCS2's schedule a fixed lead after the decision.
  \hot~trained; \frz~frozen.}
\label{fig:overview}
\end{figure*}


Legged robots can cross broken or narrow terrain that stops wheels
and tracks, as long as they find where to place their feet.
Choosing those footholds is only part of the problem: the robot must
also decide which legs to move, when to lift them, and how long its swings
should last. Together, these decisions are called
\emph{contact planning}. Most controllers simplify contact planning
with a fixed gait, which sets the leg timing and ordering in advance but cannot
react when safe footholds appear at irregular places and times.
Learned footstep planners can make these choices implicitly, and
simulation lets them practice on large amounts of terrain. Yet a learned planner is trained on top of one low-level controller and
then deployed on another, and a controller does not always place the
foot where it is told. Our \emph{training controller}, for example,
lands the foot a few centimeters off of the commanded target, which is more than the safety margin around terrain elements.
\emph{OCS2} \cite{ocs2}, a nonlinear model predictive controller
(NMPC) with whole-body control (WBC), misses in a different way and
also starts each step on its own schedule. When a learned planner fails
on a new controller, the fault therefore need not lie in choosing a
safe foothold; it can lie in landing on it.

Model-based planners place
foothold choice, contact timing, and whole-body motion inside one
model predictive controller \cite{Wensing2022Nov}. This works well
when the contact pattern is known, but once the planner must choose
contacts, the problem becomes a mixed-integer program. Some methods
reduce the search by limiting footholds to convex terrain regions
\cite{deits_footstep_2014}, and others optimize contact schedules and
timing together \cite{winkler_gait_2018}. Gait-free planners remove the
periodic schedule altogether, choosing contact timing online from
% placeholder image
\begin{figure}[H]
    \centering
    \includegraphics[width=0.5\textwidth]{example-image-a}
    \caption{Robot navigating challenging terrain.}
    \label{fig:placeholder}
\end{figure}
per-leg feasibility rules \cite{bjelonic_whole-body_2021} or by tree
search \cite{amatucci_monte_2022}, and recent work speeds up this
search by learning its cost-to-go \cite{taouil_non-gaited_2024} or by
learning per-foot flight phases inside an MPC horizon
\cite{patrizi_rl-augmented_2026}. These planners decide contact
timing, but not footholds chosen from a perceived map of safe cells.

Learned planners make the contact decision instead. End-to-end
policies put it inside one neural controller and cross hard terrain
reliably \cite{lee_learning_2020, zhang_learning_2023}; trained
without an explicit schedule, they can even develop gaits on their own
\cite{fu_minimizing_2021}. Their contact decision, however, is hidden
inside the joint targets, so it cannot be checked before execution or
handed to a model-based controller. Perceptive model-based stacks
instead project an explicit foothold onto a convex terrain region
\cite{grandia_perceptive_2022}, and hybrid methods learn that foothold
while a gait generator still fixes the timing
\cite{gangapurwala_rloc_2022, tsounis_deepgait_2020}. The closest
work, ContactNet \cite{bratta_contactnet_2024}, ranks discretized
foothold candidates with a network without prescribing which leg
swings. Its concurrency, however, is set by configuration, it re-ranks
candidates only at touchdown, and its gait sets the swing timing.

Finally, work on moving learned policies to new robots or domains has
moved body motion, joint actions, or footholds within a fixed gait,
but not gait-free footstep choices. Li et al.
\cite{li_zero-shot_2022} retarget learned quadruped locomotion
policies to new robots through a hybrid kinodynamic MPC, and RLOC
\cite{gangapurwala_rloc_2022} moves its learned footstep policy to a
larger robot while keeping its gait. Xie et al. \cite{xie_fada_2026}
split a humanoid policy into a planner, which states the intended
motion, and an inverse dynamics model, which turns it into actions,
and then adapt only the inverse model to a new domain.

We present GaitNet (Fig.~\ref{fig:overview}), which separates choosing
a step from executing it. Every \emph{tick}, the planner picks one step
or a \emph{no-step option} from a few safe \emph{candidate} footholds
per leg, kept away from \emph{voids}, the holes in the ground. Because
a controller may not place the foot where it is told, an execution
layer then makes the controller land each step where and when the
planner meant. On the training controller, this layer is an
\emph{inverse footstep model} (IDM), fitted from the controller's own
logged steps, which turns where GaitNet wants the foot into the command
that lands it there. For a new controller, the IDM is refit on that
controller's own logged steps, and on OCS2 steps are also delivered on
OCS2's own schedule. GaitNet is trained in Isaac Lab
\cite{mittal_orbit_2023} and is never retrained. It builds on the
one-step-per-tick action space of \cite{Sullivan2025}, which was tested
on one controller; the candidate attention, the motion history, the
execution layer, the second controller, and all results are new here,
and all results are in simulation.

% Kept in one block so the list never breaks across columns or pages.
% NOTE (owen): we'll want to re-condiser this later. as of now, it left a huge gap on pg2
% \noindent\begin{minipage}{\linewidth}
Our contributions are:
\begin{itemize}
  \item GaitNet, a footstep planner that is given no gait. Every tick, it
    picks one step or none from sampled safe footholds, comparing the
    candidates with attention and using a short motion history
    (Secs.~\ref{subsec:method-step} and~\ref{subsec:method-network}).
  \item An execution layer that lets the frozen GaitNet run on a new
    controller, with an inverse footstep model per controller fitted by
    hindsight relabelling from that controller's own logged steps
    (Secs.~\ref{subsec:method-idm} and~\ref{subsec:method-deployment}).
  \item An evaluation on OCS2 with a simulated depth camera, over
    0--\corlOcsMaxVoidPct\% voids and
    \corlOcsSpeedMin--\corlOcsSpeedMax\,m/s, in which GaitNet matches
    OCS2's own trot on sparse voids and exceeds it on dense voids and
    at faster commands, and ablations of overlapping swings, the IDM,
    and the candidate comparison (Sec.~\ref{sec:experiments}).
\end{itemize}
% \end{minipage}
